\chapter{Partial Equilibrium and Welfare}\label{ch:b1:partial-equilibrium-welfare}

On 1 July 2020 Germany cut its standard rate of value-added tax from $19\%$
to $16\%$ and its reduced rate, which covers most food, from $7\%$ to $5\%$,
for six months. % ledger: b1:de-vat-2020
Consumer-price inflation, measured over twelve months, fell from $0.8\%$ in
June to $0.0\%$ in July; food inflation fell from $4.1\%$ to $0.9\%$. When
the rates went back up in January 2021, inflation jumped from $-0.7\%$ in
December to $1.6\%$. % figdata: de_hicp_points
A tax cut changed prices within a month, and by less than the cut itself.
Who gets a tax cut, who pays a tax, and what does a tax cost beyond the
revenue it raises? This chapter answers with the tools of the earlier
chapters: demand from the consumer's choice, supply from the firm's, and
surplus as an integral.

\begin{openingchart}{Germany: consumer-price inflation over twelve months,
month by month, January 2019 -- December 2022, all items and food; the
months of the VAT cut shaded.}{Eurostat, Harmonised index of consumer prices,
monthly annual rates of change (prc\_hicp\_manr)}
\begin{tikzpicture}
\begin{axis}[omtime, xlabel={Year}, ylabel={Inflation over 12 months (\%)}, xmin=2019, xmax=2023,
  ymin=-2, ymax=22, width=0.74\linewidth, height=0.44\linewidth, legend pos=north west,
  xtick={2019,2020,2021,2022,2023}]
  \fill[omDef!14] (axis cs:2020.5,-2) rectangle (axis cs:2021,22);
  \draw[gray, dotted] (axis cs:2019,0) -- (axis cs:2023,0);
  \addplot[demand, mark=none] table[x=x, y=all, col sep=comma] {figdata/bachelor-1/13-partial-equilibrium-welfare/out/de_hicp.csv};
  \addplot[supply, mark=none] table[x=x, y=food, col sep=comma] {figdata/bachelor-1/13-partial-equilibrium-welfare/out/de_hicp.csv};
  \legend{All items, Food and non-alcoholic beverages}
\end{axis}
\end{tikzpicture}
\end{openingchart}

The chart shows the cut twice: as a fall in the twelve-month rate while it
lasted, and as a rise in July--December 2021, when the months of the cut
became the base of the comparison and prices were compared with the low
prices of a year before. In July 2021 inflation was $3.1\%$, and food inflation
$4.3\%$. % figdata: de_hicp_points
Later the chart is dominated by the energy shock of 2022, which belongs to
another chapter.

\begin{recall}
Consumer and \omterm{def:b1:partial-equilibrium-welfare:surplus}{producer surplus}, the \omterm{def:b1:partial-equilibrium-welfare:dwl}{deadweight loss} of a tax and its
incidence were met at school. The incidence of a per-unit tax with general
supply and demand (\cref{prop:b1:comparative-statics:incidence}) and the
consumer's surplus as an area under demand
(\cref{prop:b1:consumer-welfare:areas}) are earlier in this book; the
integral and its fundamental theorem are from the mathematics series.
\end{recall}

\section{Equilibrium of one market}\label{sec:b1:partial-equilibrium-welfare:equilibrium}

\begin{definition}[Partial equilibrium]\label{def:b1:partial-equilibrium-welfare:partial}
\index{partial equilibrium}
A \emph{partial \omterm{def:b1:models-and-data:equilibrium}{equilibrium}} is a price of one good at which the quantity
demanded equals the quantity supplied, with the prices of all other goods,
and incomes, held fixed.
\end{definition}

Write $Z(p) = D(p) - S(p)$, the excess demand. If $D$ and $S$ are continuous, if
$Z(p) > 0$ at some low price and $Z(p) < 0$ at some high price, then by the
intermediate value theorem there is a price between them with $Z(p) = 0$: an
\omterm{def:b1:models-and-data:equilibrium}{equilibrium} exists. If $D$ is decreasing and $S$ increasing, $Z$ is strictly
decreasing and the \omterm{def:b1:models-and-data:equilibrium}{equilibrium} is unique.

\omterm{def:b1:models-and-data:equilibrium}{Equilibrium} is a state, not a process. A simple account of how a market
gets there is that the price rises when there is excess demand and falls
when there is excess supply: $\dot p = k\,Z(p)$, with $k > 0$.

\begin{definition}[Stable equilibrium]\label{def:b1:partial-equilibrium-welfare:stable}
\index{stable equilibrium}\index{stability (of an equilibrium)}
An \omterm{def:b1:models-and-data:equilibrium}{equilibrium} price $p^*$ is \emph{stable} under the price adjustment
$\dot p = kZ(p)$ if the price returns to $p^*$ from any price close enough to
it.
\end{definition}

\begin{proposition}[Stability]\label{prop:b1:partial-equilibrium-welfare:stability}
If $D'(p^*) < S'(p^*)$, the \omterm{def:b1:models-and-data:equilibrium}{equilibrium} $p^*$ is \omterm{def:b1:partial-equilibrium-welfare:stable}{stable}; if $D'(p^*) > S'(p^*)$, it is
unstable.
\end{proposition}

\begin{proof}
Near $p^*$, $Z(p) \approx (D'(p^*) - S'(p^*))(p - p^*)$, so $x = p - p^*$ follows
$\dot x = k(D' - S')x$, whose solution is $x(t) = x(0)e^{k(D' - S')t}$. It tends to zero
if $D' - S' < 0$ and grows if $D' - S' > 0$.
\end{proof}

With downward-sloping demand and upward-sloping supply the condition
always holds. It fails only when quantity demanded rises with the price
faster than quantity supplied ($D' > S'$) --- with a \omterm{def:b1:income-and-substitution-effects:giffen}{Giffen good}, say --- or
quantity supplied falls with the price faster than quantity demanded. In farm markets, where output is decided a season before it is
sold, prices can overshoot instead: a high price this year brings a large
crop and a low price next year, which brings a small crop and a high
price, and so on, a cycle called the \emph{cobweb}, which dies out when
supply responds less to price than demand does.

\section{Surplus as an integral}\label{sec:b1:partial-equilibrium-welfare:surplus}

\begin{definition}[Producer and total surplus]\label{def:b1:partial-equilibrium-welfare:surplus}
\index{producer surplus}\index{total surplus}\index{surplus!producer}\index{surplus!total}
At the price $p^*$, \emph{producer surplus} is the area between the price and
the supply curve, $\mathrm{PS} = \int_0^{p^*}S(p)\,\dd p$: revenue minus the variable
cost of the units sold. \emph{Total surplus} is \omterm{def:b1:consumer-welfare:consumer-surplus}{consumer surplus},
$\mathrm{CS} = \int_{p^*}^{\bar p}D(p)\,\dd p$ (with $\bar p$ the price at which demand
reaches zero), plus producer surplus, plus any tax revenue.
\end{definition}

In terms of quantity, with $P^d(Q)$ the inverse demand (the most a buyer will
pay for the $Q$-th unit) and $P^s(Q)$ the inverse supply (the marginal cost of
the $Q$-th unit), \omterm{def:b1:partial-equilibrium-welfare:surplus}{total surplus} from selling $Q$ units is
$W(Q) = \int_0^Q\bigl(P^d(x) - P^s(x)\bigr)\,\dd x$.

\begin{proposition}[Efficiency of competitive equilibrium]\label{prop:b1:partial-equilibrium-welfare:efficiency}
If $P^d$ is decreasing and $P^s$ increasing, $W(Q)$ is largest at the
competitive quantity, where $P^d(Q) = P^s(Q)$.
\end{proposition}

\begin{proof}
By the fundamental theorem of calculus, $W'(Q) = P^d(Q) - P^s(Q)$, which is
positive below the competitive quantity and negative above it.
\end{proof}

The proposition says that every unit worth more to some buyer than its
cost to some seller is traded, and no other. It says nothing about who
gets the surplus, and it assumes that the curves measure all the benefits
and costs: no one else is affected by the trade.

\begin{omfigure}
\begin{tikzpicture}
\begin{axis}[omrecord, grid=none, xlabel={Quantity $Q$}, ylabel={Price $p$}, xmin=0, xmax=50, ymin=0, ymax=7.5,
  width=0.52\linewidth, height=0.4\linewidth, samples=100]
  \addplot[cs area, draw=none, domain=0:32, forget plot] {sqrt(48 - x)} -- (axis cs:32,4) -- (axis cs:0,4) -- cycle;
  \addplot[fill=omThm!20, draw=none, domain=0:32, forget plot] {sqrt(x/2)} -- (axis cs:32,4) -- (axis cs:0,4) -- cycle;
  \addplot[demand, domain=0:48] {sqrt(48 - x)};
  \addplot[supply, domain=0:50] {sqrt(x/2)};
  \addplot[guide] coordinates {(0,4) (32,4) (32,0)};
  \node[font=\footnotesize] at (axis cs:8,5.1) {CS};
  \node[font=\footnotesize] at (axis cs:16,3.2) {PS};
\end{axis}
\end{tikzpicture}

{\small Demand $D(p) = 48 - p^2$ (blue) and supply $S(p) = 2p^2$ (red) meet at $p = 4$,
$Q = 32$. \omterm{def:b1:consumer-welfare:consumer-surplus}{Consumer surplus} is $\int_4^{\sqrt{48}}(48 - p^2)\,\dd p = 51.03$; \omterm{def:b1:partial-equilibrium-welfare:surplus}{producer surplus} is
$\int_0^4 2p^2\,\dd p = 42.67$.}
\end{omfigure}

\section{Deadweight loss}\label{sec:b1:partial-equilibrium-welfare:dwl}

A tax of $\tau$ per unit drives a wedge between the price buyers pay, $P$, and
the price sellers keep, $p = P - \tau$. The quantity falls from $Q_0$ to $Q_\tau$,
where $P^d(Q_\tau) - P^s(Q_\tau) = \tau$.

\begin{definition}[Deadweight loss]\label{def:b1:partial-equilibrium-welfare:dwl}
\index{deadweight loss}\index{Harberger triangle}
The \emph{deadweight loss} of a tax is the fall in \omterm{def:b1:partial-equilibrium-welfare:surplus}{total surplus}, revenue
included:
\[
\mathrm{DWL}(\tau) = \int_{Q_\tau}^{Q_0}\bigl(P^d(x) - P^s(x)\bigr)\,\dd x,
\]
the surplus of the units no longer traded. Drawn with straight lines it is
a triangle, the \emph{Harberger triangle}.
\end{definition}

\begin{proposition}[The second-order formula]\label{prop:b1:partial-equilibrium-welfare:harberger}
For a small tax,
\[
\mathrm{DWL} \approx \tfrac12\tau^2\left|\frac{\dd Q}{\dd\tau}\right| = \tfrac12\tau^2\,\frac Qp\,\frac{\varepsilon_s|\varepsilon_d|}{\varepsilon_s + |\varepsilon_d|},
\]
with the elasticities and $p$, $Q$ taken at the untaxed \omterm{def:b1:models-and-data:equilibrium}{equilibrium}. The loss
grows with the square of the tax.
\end{proposition}

\begin{proof}
The integrand rises from 0 at $Q_0$ to $\tau$ at $Q_\tau$; to first order in $\tau$ it is
linear, so the integral is the triangle $\tfrac12\tau(Q_0 - Q_\tau)$, and
$Q_0 - Q_\tau \approx \tau|\dd Q/\dd\tau|$. From
\cref{prop:b1:comparative-statics:incidence},
$\dd Q/\dd\tau = D'S'/(S' - D')$; multiply numerator and denominator by $p/Q$
twice to get the elasticities.
\end{proof}

The square rule has a practical consequence: two taxes of 5 cost less than
one of 10 (since $2 \times 25 < 100$), so a given revenue is raised more cheaply
from a broad base at a low rate than from a narrow base at a high rate. The
loss also grows with the elasticities: taxing goods with elastic demand or
supply costs most.

\begin{omfigure}
\begin{tikzpicture}
\begin{groupplot}[group style={group size=2 by 1, horizontal sep=18mm}, omrecord, grid=none,
  width=0.45\linewidth, height=0.42\linewidth, samples=60]
\nextgroupplot[xlabel={Sugar $Q$ (thousand tonnes)}, ylabel={Dalens a kilogram}, xmin=0, xmax=120, ymin=0, ymax=12,
  title={\small A tax of 1.2}, width=0.45\linewidth]
  \fill[omDef!20] (axis cs:0,5.4) rectangle (axis cs:54,6.6);
  \fill[omProp!35] (axis cs:54,6.6) -- (axis cs:60,6) -- (axis cs:54,5.4) -- cycle;
  \addplot[demand, domain=0:120] {12 - x/10};
  \addplot[supply, domain=0:120] {x/10};
  \addplot[guide] coordinates {(0,6.6) (54,6.6) (54,0)};
  \addplot[guide] coordinates {(0,5.4) (54,5.4)};
  \node[font=\footnotesize, anchor=west] at (axis cs:62,6) {DWL};
  \node[font=\footnotesize] at (axis cs:27,6) {revenue};
\nextgroupplot[xlabel={Tax $\tau$ (dalens a kilogram)}, ylabel={Millions of dalens a year}, xmin=0, xmax=12,
  ymin=0, ymax=200, title={\small Revenue and loss}, width=0.45\linewidth, legend pos=north west,
  legend style={fill=none, draw=none}]
  \addplot[demand, domain=0:12] {x*(60 - 5*x)};
  \addplot[supply, domain=0:8.5] {2.5*x^2};
  \legend{Revenue, DWL}
\end{groupplot}
\end{tikzpicture}

{\small Left: Brenland's sugar market (invented), demand $Q = 120 - 10P$ and supply
$Q = 10p$, with a tax of 1.2 dalens a kilogram: buyers pay 6.6, sellers keep
5.4, revenue is 64.8 and the \omterm{def:b1:partial-equilibrium-welfare:dwl}{deadweight loss} 3.6 (millions of dalens a year).
Right: as the tax rises, revenue $\tau(60 - 5\tau)$ peaks at $\tau = 6$ and then falls,
while the loss $2.5\tau^2$ keeps rising.}
\end{omfigure}

\section{Incidence}\label{sec:b1:partial-equilibrium-welfare:incidence}

\begin{definition}[Statutory and economic incidence]\label{def:b1:partial-equilibrium-welfare:incidence}
\index{statutory incidence}\index{economic incidence}\index{tax incidence}
The \emph{statutory incidence} of a tax says who must pay it to the
government; the \emph{economic incidence} says whose surplus it reduces,
through the changes in the prices they pay and receive.
\end{definition}

In a competitive market the two need not be related: whether sellers or
buyers hand over the tax, the wedge between $P$ and $p$ is the same, and so is
everyone's surplus. Buyers bear the share $\varepsilon_s/(\varepsilon_s + |\varepsilon_d|)$ of a small
tax and sellers the rest (\cref{prop:b1:comparative-statics:incidence}):
the side that can less easily walk away pays more.

\begin{definition}[Pass-through rate]\label{def:b1:partial-equilibrium-welfare:pass-through}
\index{pass-through rate}
The \emph{pass-through rate} of a tax is the change in the price buyers pay
per unit change in the tax, $\dd P/\dd\tau$; for a percentage tax, the
percentage change in the buyers' price as a share of the percentage change
in $1 + t$.
\end{definition}

A competitive market with a fixed number of firms has a \omterm{def:b1:partial-equilibrium-welfare:pass-through}{pass-through rate}
between 0 and 1. A rate above one is possible with market power
(\cref{ch:b1:monopoly-pricing}); in the long run of a \omterm{def:b1:competitive-industry:cost-industry}{constant-cost
industry} it is exactly one (\cref{ch:b1:competitive-industry}).

\begin{record}[A temporary VAT cut]
The German cut of 2020 lowered the gross price of a good taxed at the
standard rate by $1 - 1.16/1.19 = 2.5\%$ and of one at the reduced rate by
$1 - 1.05/1.07 = 1.9\%$, if fully passed on. A study of daily prices of about
130,000 supermarket products, compared with the same products in Austria,
found that prices fell by about $1.3\%$: some $70\%$ of the cut was passed on
to buyers, more where many competing products were sold. The return to the
old rates raised prices by only about half as much as the cut had lowered
them (Fuest, Neumeier and Stöhlker 2021). % ledger: b1:fuest2021
The competitive model predicts the partial pass-through; it does not
predict the asymmetry, which suggests that sellers set prices and adjust
them at a cost.
\end{record}

\begin{measure}[Monthly inflation from the HICP]
The harmonised index of consumer prices measures the price of a fixed
basket in each month; the rate in the opening chart compares each month
with the same month a year before. A one-off fall in the price level
therefore lowers the twelve-month rate for twelve months and then drops out
of it. A temporary cut that is reversed six months later, like the German
one, lowers the rate while it lasts and raises it during the six months in
which the low months serve as the base. Neither changes inflation for good.
\end{measure}

\section{When one market can be studied alone}\label{sec:b1:partial-equilibrium-welfare:alone}

Everything above held other prices and incomes fixed. That is a good
approximation when the good takes a small share of buyers' budgets (so
\omterm{def:b1:income-and-substitution-effects:effects}{income effects} are small and the CV, EV and the \omterm{def:b1:consumer-welfare:consumer-surplus}{consumer surplus} agree,
\cref{ch:b1:consumer-welfare}), when its market is small relative to the
markets for its inputs, and when it has no close substitutes or
complements whose prices would move. When a taxed good has a close
substitute, demand for the substitute rises, its price rises if its supply
slopes up, and surplus changes in that market too. If the substitute's
market has no tax of its own, those changes are transfers between its
buyers and sellers, and the loss is the triangle under a demand curve for
the taxed good that lets the substitute's price adjust; the triangle drawn
with the substitute's price held fixed overstates it
(\cref{exo:b1:partial-equilibrium-welfare:substitute}). Taxes on large
markets --- labour, housing --- move other prices and need the general
\omterm{def:b1:models-and-data:equilibrium}{equilibrium} of Year~2.

\begin{model}{Competitive partial equilibrium}
\assumes One market, price-taking buyers and sellers, other prices and
incomes fixed, no effects on third parties.
\predicts Buyers bear the share $\varepsilon_s/(\varepsilon_s + |\varepsilon_d|)$ of a tax whoever pays it to
the state; a \omterm{def:b1:partial-equilibrium-welfare:dwl}{deadweight loss} that grows with the square of the tax.
\tested About $70\%$ of Germany's 2020 VAT cut reached supermarket prices
(Fuest, Neumeier and Stöhlker 2021); % ledger: b1:fuest2021
cigarette taxes are less than fully passed through, and less where buyers
can cross to a state with lower taxes (Harding, Leibtag and Lovenheim
2012). % ledger: b1:harding2012
\breaks With market power; with temporary changes, where the cut and the
return were passed through unequally; for large markets whose prices move
others.
\end{model}

\section*{Exercises}\label{sec:b1:partial-equilibrium-welfare:exercises}

\begin{exercise}[$\star$]\label{exo:b1:partial-equilibrium-welfare:cs-ps}
Find the \omterm{def:b1:models-and-data:equilibrium}{equilibrium}, \omterm{def:b1:consumer-welfare:consumer-surplus}{consumer surplus} and \omterm{def:b1:partial-equilibrium-welfare:surplus}{producer surplus} for $D = 120 - 2p$
and $S = p$.
\end{exercise}

\begin{exercise}[$\star$]\label{exo:b1:partial-equilibrium-welfare:read}
From the opening chart's table, German inflation was $0.8\%$ in June 2020,
$0.0\%$ in July, $-0.7\%$ in December and $1.6\%$ in January 2021; food inflation
was $4.1\%$ in June and $0.9\%$ in July. % figdata: de_hicp_points
Compute the changes in points at the start and at the end of the cut.
\end{exercise}

\begin{exercise}[$\star$]\label{exo:b1:partial-equilibrium-welfare:share}
With $\varepsilon_d = -0.5$ and $\varepsilon_s = 1.5$, what share of a small tax do buyers bear?
\end{exercise}

\begin{exercise}[$\star$]\label{exo:b1:partial-equilibrium-welfare:cigarettes}
Smokers' demand for cigarettes is inelastic and tobacco firms can sell in
many markets. Predict which side bears more of a cigarette tax.
\end{exercise}

\begin{exercise}[$\star\star$]\label{exo:b1:partial-equilibrium-welfare:cs-integral}
Compute the loss of \omterm{def:b1:consumer-welfare:consumer-surplus}{consumer surplus} when the price of a good with demand
$D(p) = 100/p$ rises from 1 to 2.
\end{exercise}

\begin{exercise}[$\star\star$]\label{exo:b1:partial-equilibrium-welfare:exact}
Demand is $D(p) = 100/p$ and supply $S(p) = p$. A tax of $10\%$ of the sellers'
price is imposed. Compute the prices and the exact \omterm{def:b1:partial-equilibrium-welfare:dwl}{deadweight loss}, and
compare with the second-order formula evaluated with a tax of one dalen
(the tax at the untaxed price) and with the actual wedge.
\end{exercise}

\begin{exercise}[$\star\star$]\label{exo:b1:partial-equilibrium-welfare:double}
A tax is doubled. Predict, to second order, what happens to the \omterm{def:b1:partial-equilibrium-welfare:dwl}{deadweight
loss}, and why revenue less than doubles.
\end{exercise}

\begin{exercise}[$\star\star$]\label{exo:b1:partial-equilibrium-welfare:stability-ex}
With $\dot p = 0.5\,(D(p) - S(p))$ and $S(p) = 2p$, study the stability of the
\omterm{def:b1:models-and-data:equilibrium}{equilibrium} for $D(p) = 20 + p$ and for $D(p) = -20 + 3p$.
\end{exercise}

\begin{exercise}[$\star\star$]\label{exo:b1:partial-equilibrium-welfare:vat}
Compute the fall in gross prices that full pass-through of the 2020 cut
predicts at the standard and the reduced rate. Compare with the fall in
the twelve-month inflation rates of the opening chart's table: why is the
comparison only rough?
\end{exercise}

\begin{exercise}[$\star\star\star$]\label{exo:b1:partial-equilibrium-welfare:max}
Prove that \omterm{def:b1:partial-equilibrium-welfare:surplus}{total surplus} is largest at the competitive quantity when the
inverse demand is decreasing and the inverse supply increasing, and show
that a quantity above it also loses surplus.
\end{exercise}

\begin{exercise}[$\star\star\star$]\label{exo:b1:partial-equilibrium-welfare:substitute}
Butter is taxed. Margarine, a close substitute, has an upward-sloping
supply. Show that margarine's producers gain surplus and its buyers lose
it, and explain why the \omterm{def:b1:partial-equilibrium-welfare:dwl}{deadweight loss} measured in the butter market
alone, with the margarine price held fixed, is not the whole loss.
\end{exercise}

\begin{exercise}[$\star\star\star$]\label{exo:b1:partial-equilibrium-welfare:refute}
What pattern in the German prices of 2020--21 would refute competitive
incidence as an account of the VAT cut? Name one comparison across
products and one over time.
\end{exercise}

\section{Problem: A Sugar Tax in Brenland}\label{sec:b1:partial-equilibrium-welfare:problem}

\begin{problem}[{Weekend problem --- the marginal deadweight loss of a tax}]\label{pb:b1:partial-equilibrium-welfare:sugar}
Brenland's demand for sugar (invented) is $Q = 120 - 10P$ and its supply is
$Q = 10p$, in thousands of tonnes a year, where $P$ is the price buyers pay and
$p$ the price sellers keep, in dalens a kilogram. A price in dalens a
kilogram times a quantity in thousands of tonnes is a sum in millions of
dalens: all surpluses and revenues below are in millions of dalens a year.

\medskip\noindent\textbf{Part I --- Equilibrium.}
\begin{enumerate}
  \item Find the untaxed \omterm{def:b1:models-and-data:equilibrium}{equilibrium}.
  \item Compute the price elasticities of demand and supply there.
  \item Write the price adjustment $\dot p = k(D - S)$ and show that the \omterm{def:b1:models-and-data:equilibrium}{equilibrium} is \omterm{def:b1:partial-equilibrium-welfare:stable}{stable}.
  \item Compute \omterm{def:b1:consumer-welfare:consumer-surplus}{consumer surplus}.
  \item Compute \omterm{def:b1:partial-equilibrium-welfare:surplus}{producer surplus}.
  \item Show that \omterm{def:b1:partial-equilibrium-welfare:surplus}{total surplus} would be smaller at $Q = 50$ and at $Q = 70$.
\end{enumerate}

\medskip\noindent\textbf{Part II --- A tax.} A tax of 1.2 dalens a kilogram ($20\%$
of the untaxed price) is collected from sellers.
\begin{enumerate}[resume]
  \item Find $P$, $p$ and the quantity.
  \item Compute the buyers' share from the elasticities and compare.
  \item Compute the revenue.
  \item Would anything change if the buyers paid the tax?
  \item What is the \omterm{def:b1:partial-equilibrium-welfare:pass-through}{pass-through rate}?
  \item Instead, the government pays a subsidy of 1.2 a kilogram. Find the quantity and the cost of the subsidy.
\end{enumerate}

\medskip\noindent\textbf{Part III --- Surplus.}
\begin{enumerate}[resume]
  \item Compute \omterm{def:b1:consumer-welfare:consumer-surplus}{consumer surplus} with the tax.
  \item Compute \omterm{def:b1:partial-equilibrium-welfare:surplus}{producer surplus} with the tax.
  \item Compute the \omterm{def:b1:partial-equilibrium-welfare:dwl}{deadweight loss} as the fall in \omterm{def:b1:partial-equilibrium-welfare:surplus}{total surplus}.
  \item Compute it with the second-order formula.
  \item Why do the two agree exactly here?
  \item The tax doubles to 2.4. Compute revenue and \omterm{def:b1:partial-equilibrium-welfare:dwl}{deadweight loss}.
  \item Compute the \omterm{def:b1:partial-equilibrium-welfare:dwl}{deadweight loss} per dalen of revenue at both rates.
\end{enumerate}

\medskip\noindent\textbf{Part IV --- How much to tax.}
\begin{enumerate}[resume]
  \item Write revenue as a function of $\tau$.
  \item Find the tax that raises the most revenue, and the revenue.
  \item Compute the \omterm{def:b1:partial-equilibrium-welfare:dwl}{deadweight loss} there, per dalen raised.
  \item Show that the \omterm{def:b1:partial-equilibrium-welfare:dwl}{deadweight loss} of the last dalen raised is $5\tau/(60 - 10\tau)$.
  \item Evaluate it at $\tau = 1.2$ and at $\tau = 2.4$.
  \item What happens to it as $\tau$ approaches 6?
  \item State the result: the marginal \omterm{def:b1:partial-equilibrium-welfare:dwl}{deadweight loss} of the last dalen raised at a $20\%$ rate.
\end{enumerate}
\end{problem}
